\begin{experiment}
\textbf{Objective}
\begin{itemize}
  \item Use readelf to examine the ELF header, sections, segments, symbols, and relocations.
\end{itemize}

\textbf{Setup requirements}
\begin{itemize}
  \item gcc or clang installed.
  \item binutils tools (readelf).
\end{itemize}

\textbf{Step-by-step commands}
\begin{enumerate}
  \item Build a small executable:
\begin{lstlisting}[language=bash]
gcc -o hello hello.c
\end{lstlisting}
  \item Inspect the ELF header:
\begin{lstlisting}[language=bash]
readelf -h hello
\end{lstlisting}
  \item Inspect section headers:
\begin{lstlisting}[language=bash]
readelf -S hello
\end{lstlisting}
  \item Inspect program headers:
\begin{lstlisting}[language=bash]
readelf -l hello
\end{lstlisting}
  \item Inspect symbols and relocations:
\begin{lstlisting}[language=bash]
readelf -s -r hello
\end{lstlisting}
\end{enumerate}

\textbf{Expected output explanation}
\begin{itemize}
  \item The ELF header identifies class, endianness, and ABI.
  \item Section headers show how data is organized.
  \item Program headers show loadable segments.
  \item Symbol and relocation tables show linkage metadata.
\end{itemize}

\textbf{Questions for reader}
\begin{enumerate}
  \item Which sections are mapped into memory at runtime?
  \item Why do some binaries have no relocation entries?
  \item How does a PIE binary change the program headers?
\end{enumerate}
\end{experiment}
